\documentclass[11pt]{article}
\linespread{1.25}
\usepackage[a4paper,left=2.5cm,right=2.5cm,top=2cm,bottom=2cm]{geometry}
\setlength{\parindent}{0pt}
\setlength{\parskip}{6pt plus 2pt minus 1pt}
\usepackage{lastpage}
\usepackage{multicol}
\pagenumbering{gobble}

% Date stuff:
\usepackage[yyyymmdd]{datetime}
\renewcommand{\dateseparator}{-}

\begin{document}
\begin{multicols}{2}
\begin{flushleft}
    \small
    Christian Vedel \\
    christian-vs@sam.sdu.dk \\
\end{flushleft}
\columnbreak
\begin{flushright}
    \small
    \today
\end{flushright}
\end{multicols}

\vspace{0.25cm}
    \Large
    \textbf{A Perfect Storm: First-Nature Geography and Economic Development}


\vspace{1cm}

Dear Editors, 

I am pleased to submit my manuscript entitled "A Perfect Storm: First-Nature Geography and Economic Development" for consideration at the American Economic Review. This paper exploits a rare and exogenous natural experiment—a storm-induced channel breach in 1825 Denmark—to estimate the long-run causal impact of first-nature geography on regional economic development.

The findings contribute to core debates in development and urban economics, showing how sudden geographic shifts shape the location of prosperity. This also informs contemporary policy discussions on climate adaptation and infrastructure planning.

Thank you for your consideration.

\vspace{0.5cm}

Sincerely,

Christian Vedel,\\
University of Southern Denmark
\end{document}
